\begin{thebibliography}{99}
\bibitem{Chen2022} 陈怡然，王一土，神经网络加速器架构概述，中国科学：信息科学，52(4)：596--611，2022。DOI: 10.1360/SSI-2021-0409。
\bibitem{Problem2026} 中国研究生数学建模竞赛组委会，通用神经网络处理器下的多核调度问题及配套附件，第二十三届中国研究生数学建模竞赛 A 题，2026。
\bibitem{Liu2025} 刘正煜，张帆，祁晓峰，高彦钊，宋怡景，范旺，深度学习编译器研究综述，计算机科学，52(8)：29--44，2025。DOI: 10.11896/jsjkx.250100062。
\bibitem{Graham1969} Graham R L，Bounds on Multiprocessing Timing Anomalies，SIAM Journal on Applied Mathematics，17(2)：416--429，1969。DOI: 10.1137/0117039。
\bibitem{Topcuoglu2002} Topcuoglu H，Hariri S，Wu M Y，Performance-effective and low-complexity task scheduling for heterogeneous computing，IEEE Transactions on Parallel and Distributed Systems，13(3)：260--274，2002。DOI: 10.1109/71.993206。
\bibitem{Yang2025} 杨紫超，吴恒，吴悦文，张文博，基于性能建模的深度学习训练任务调度综述，软件学报，36(4)：1570--1589，2025。DOI: 10.13328/j.cnki.jos.007202。
\bibitem{Chen2024} Chen R，Ding Z，Zheng S，等，MAGIS: Memory Optimization via Coordinated Graph Transformation and Scheduling for DNN，Proceedings of ASPLOS 2024，Volume 3：607--621，2024。DOI: 10.1145/3620666.3651330。
\bibitem{Parashar2019} Parashar A，Raina P，Shao Y S，等，Timeloop: A Systematic Approach to DNN Accelerator Evaluation，Proceedings of ISPASS 2019：304--315，2019。DOI: 10.1109/ISPASS.2019.00042。
\bibitem{AIGPT6} OpenAI，GPT-6（模型标识：gpt-6-astra），发布日期：2026 年 9 月 3 日。\url{https://openai.com/index/gpt-6-astra/}。
\end{thebibliography}
